\documentclass[11pt]{article}
\usepackage[margin=1in]{geometry}
\usepackage{amsmath,amssymb,amsthm}
\usepackage[colorlinks=true,urlcolor=blue,linkcolor=blue]{hyperref}
\newtheorem{theorem}{Theorem}
\newtheorem{lemma}{Lemma}
\newtheorem{remark}{Remark}
\newcommand{\tr}{\operatorname{tr}}
\newcommand{\rank}{\operatorname{rank}}
\newcommand{\N}{\mathcal N}
\title{A uniform relative compression bound on unit-weight paths}
\author{Independent path upper-bound calculation}
\date{November 15, 2026}
\begin{document}
\maketitle

Let $L_n$ be the ordinary combinatorial Laplacian of the path on
$n\ge2$ vertices, all edge conductances and vertex masses equal to one.
For $1\le p<n$, set $b=e_p-e_{p+1}$ and $L'_n=L_n+bb^T$; thus precisely
one edge conductance is doubled. Fix $0<\alpha<1$, and write
\[
 C_{n,p}=(L_n^\dagger)^{\alpha/2}
       ((L'_n)^\alpha-L_n^\alpha)(L_n^\dagger)^{\alpha/2},\qquad
 \N_{n,p}(\varepsilon)=\#\{j:\lambda_j(C_{n,p})>\varepsilon\}.
\]
Every inverse power is zero on the constant vectors. All calculations
below may therefore be carried out on their orthogonal complement.

\begin{theorem}[Uniform path upper bound]\label{thm:path}
There is a constant $K_\alpha<\infty$ such that, for every $n,p$ and
$0<\varepsilon<1$,
\[
 \N_{n,p}(\varepsilon)
       \le K_\alpha(1+\log(1/\varepsilon))^2.
\]
More explicitly, with $L=\log(1/\varepsilon)$,
\[
 \sup_{n,p}\N_{n,p}(e^{-L})
 \le\left[\frac{(1+2\alpha)^2}{4\pi^2\alpha}+o_\alpha(1)\right]L^2
 \quad(L\longrightarrow\infty).
\]
The displayed leading coefficient is only an upper bound and is not
claimed to be optimal. The weighted-path
$\Theta_\alpha(1/\varepsilon)$ relative-count law therefore fails for
unit-weight paths.
\end{theorem}

The proof uses a low-energy leverage measure rather than a bound on
the condition number of $L_n$. No restriction on $n$ relative to the
accuracy is imposed.

\begin{lemma}[Relative trace on a low-energy spectral subspace]
Let $A\succ0$, $B=A+uu^T$, and define
$C=A^{-\alpha/2}(B^\alpha-A^\alpha)A^{-\alpha/2}$. If
$A\phi_k=\lambda_k\phi_k$, put $z_k=|\langle u,\phi_k\rangle|^2/\lambda_k$
and
\[
 \nu([0,\delta])=\sum_{\lambda_k\le\delta}z_k.
\]
If $Q=1_{(0,\delta]}(A)$, then
\[
 C\succeq0,\qquad \tr(QCQ)\le\alpha\nu([0,\delta]).
\]
\end{lemma}
\begin{proof}
For $c_\alpha=\sin(\pi\alpha)/\pi$, the fractional resolvent formula
and Sherman--Morrison identity give
\[
 C=c_\alpha\int_0^\infty
 \frac{t^\alpha v_t v_t^T}{1+u^T(A+tI)^{-1}u}\,dt,
 \qquad v_t=A^{-\alpha/2}(A+tI)^{-1}u.
\]
Omitting the denominator gives a positive upper bound $D\succeq C$.
Its diagonal in the eigenbasis of $A$ is
\[
 D_{kk}=c_\alpha|\langle u,\phi_k\rangle|^2\lambda_k^{-\alpha}
       \int_0^\infty\frac{t^\alpha}{(\lambda_k+t)^2}\,dt
       =\alpha z_k.
\]
The integral is the scalar derivative of $x^\alpha$ at
$x=\lambda_k$. Summing over the range of $Q$ proves the claim.
\end{proof}

\begin{lemma}[Exact path leverage measure]\label{lem:measure}
For the unit-weight path and any edge $p$, one has
\[
 \nu_{n,p}([0,\delta])\le\sqrt\delta\qquad(\delta>0).
\]
\end{lemma}
\begin{proof}
The nonconstant path eigenpairs are, for $1\le k<n$,
\[
 \lambda_k=4\sin^2\frac{k\pi}{2n},\qquad
 \phi_k(j)=\sqrt{\frac2n}\cos\frac{k\pi(j-1/2)}n.
\]
Taking the edge difference cancels the factor in $\lambda_k$ and gives
the exact identity
\[
 z_k=\frac{|\phi_k(p)-\phi_k(p+1)|^2}{\lambda_k}
       =\frac2n\sin^2\frac{k\pi p}n\le\frac2n.
\]
On $[0,\pi/2]$, $\sin x\ge2x/\pi$, so
$\lambda_k\ge4k^2/n^2$. At most $n\sqrt\delta/2$ positive path
eigenvalues can be at most $\delta$, proving the bound.
\end{proof}

\begin{lemma}[PSD pinching count bound]\label{lem:pinch}
Let $C\succeq0$, let $P,Q$ be complementary orthogonal projections,
and suppose $\|QCQ\|\le\varepsilon/2$. Then
\[
 \#\{\lambda_j(C)>\varepsilon\}
 \le\#\{\lambda_j(PCP)>\varepsilon/2\}.
\]
\end{lemma}
\begin{proof}
The positive block matrix satisfies
$C\preceq2(PCP+QCQ)$: apply
$2\langle C^{1/2}Px,C^{1/2}Qx\rangle
\le\|C^{1/2}Px\|^2+\|C^{1/2}Qx\|^2$.
The operator on the right is block diagonal. Its $Q$ block has no
eigenvalue strictly exceeding $\varepsilon$, and its $P$ block is
$2PCP$. Eigenvalue monotonicity proves the strict count statement.
\end{proof}

\begin{proof}[Proof of Theorem~\ref{thm:path}]
Put $\delta=(\varepsilon/(2\alpha))^2$ and take
$Q=1_{(0,\delta]}(L_n)$ and $P=I-Q$ on the nonconstant space.
The two measure lemmas imply
\[
 \|QCQ\|\le\tr(QCQ)\le\alpha\sqrt\delta=\varepsilon/2.
\]
Thus Lemma~\ref{lem:pinch} reduces the desired count to that of $PCP$
at threshold $\varepsilon/2$.

Both endpoint spectra lie in $[0,6]$, since $\|L_n\|\le4$ and
$\|bb^T\|=2$. Stahl's rational approximation theorem supplies a real
rational function $r_q$ of type $(q,q)$, with no poles on $[0,6]$, such
that
\[
 \sup_{x\in[0,6]}|r_q(x)-x^\alpha|
 \le K_\alpha\exp(-2\pi\sqrt{\alpha q}).
\]
Because $\rank(L'_n-L_n)=1$, the standard rational rank identity gives
\[
 \rank(r_q(L'_n)-r_q(L_n))\le q.
\]
For completeness, if $r_q=a_q/d_q$, factor
$a_q(x)d_q(y)-d_q(x)a_q(y)=(x-y)\sum_{i,j<q}c_{ij}x^iy^j$.
Left-right evaluation at $x=L'_n,y=L_n$ shows that the range of
$d_q(L'_n)[r_q(L'_n)-r_q(L_n)]d_q(L_n)$ lies in
$\operatorname{span}\{(L'_n)^ib:0\le i<q\}$.
The denominators are invertible, proving the rank bound.

Writing $\Delta=(L'_n)^\alpha-L_n^\alpha$, spectral calculus gives
\[
 \|\Delta-(r_q(L'_n)-r_q(L_n))\|
 \le2K_\alpha\exp(-2\pi\sqrt{\alpha q}).
\]
Compressing by $P(L_n^\dagger)^{\alpha/2}$ on both sides increases
this error by at most $\delta^{-\alpha}$. Therefore $PCP$ has a rank
at most $q$ approximant with error at most $\varepsilon/2$ whenever
\[
 2K_\alpha\exp(-2\pi\sqrt{\alpha q})
       \le(\varepsilon/2)\delta^\alpha.
\]
Choosing
\[
 q=\frac{[(1+2\alpha)L+O_\alpha(1)]^2}{4\pi^2\alpha}+O(1)
\]
is sufficient. The best-rank approximation theorem and the pinching
count bound prove the asymptotic statement. Enlarging a fixed constant
handles the remaining bounded accuracy range, proving the first claim.
\end{proof}

\begin{remark}[A general sufficient condition]
Exactly the same proof applies to any finite-dimensional rank-one
family with uniformly bounded endpoint spectra and with
$\nu([0,\delta])\le K\delta^\beta$ for fixed $K,\beta>0$.
It gives $O(\log^2(1/\varepsilon))$ relative count uniformly in the
family; the nonsharp coefficient from this proof is
$(1+\alpha/\beta)^2/(4\pi^2\alpha)$.
This sufficient condition does not follow merely from unit edge
conductances or bounded maximum degree on general graphs. For example,
an edge joining two large well-connected pieces can carry a fixed
amount of leverage in an eigenvalue tending to zero.
\end{remark}

\begin{thebibliography}{1}
\bibitem{Stahl}
H.~Stahl, \emph{Best uniform rational approximation of $x^\alpha$ on
$[0,1]$}, Bull. Amer. Math. Soc. (N.S.) 28 (1993), 116--122.
\url{https://arxiv.org/abs/math/9301217}.
The root-exponential estimate used here is Theorem~1, equation~(2.1).
The full primary Theorem~1 was inspected in the parent research session
on November 15, 2026.
\end{thebibliography}
\end{document}
